% from ega2-1-fr.tex
\oldpage[II]{94}

La nécessité de la condition résulte de
(\hyperref[II.4.6.13-fr]{4.6.13, (i bis)}). Pour voir que la condition est
suffisante, on peut se borner au cas où $Y$ est affine, et alors le fait
que $\mathcal{L}''$ est ample résulte du critère
(\hyperref[II.4.5.2-fr]{4.5.2, a)}) appliqué à $\mathcal{L}$,
$\mathcal{L}'$ et $\mathcal{L}''$, en observant qu'une section de
$\mathcal{L}$ au-dessus de $X$ se prolonge (par $0$) en une section de
$\mathcal{L}''$ au-dessus de $X''$.

\begin{proposition}[4.6.18]
\label{II.4.6.18-fr}
Soient $Y$ un préschéma quasi-compact, $\mathcal{S}$ une
$\mathcal{O}_Y$-Algèbre graduée quasi-cohérente de type fini,
$X=\operatorname{Proj}(\mathcal{S})$, $f:X\to Y$ le morphisme structural.
Alors $f$ est de type fini, et il existe un entier $d>0$ tel que
$\mathcal{O}_X(d)$ soit inversible et $f$-ample.
\end{proposition}

En vertu de (\hyperref[II.3.1.10-fr]{3.1.10}), il existe un entier $d>0$
tel que $\mathcal{S}^{(d)}$ soit engendrée par $\mathcal{S}_d$. On sait que,
dans l'isomorphisme canonique entre $X$ et
$X^{(d)}=\operatorname{Proj}(\mathcal{S}^{(d)})$, $\mathcal{O}_X(d)$
s'identifie à $\mathcal{O}_{X^{(d)}}(1)$
(\hyperref[II.3.2.9-fr]{3.2.9, (ii)}). On voit donc qu'on est ramené au cas
où $\mathcal{S}$ est engendrée par $\mathcal{S}_1$~; la proposition résulte
alors de (\hyperref[II.4.4.3-fr]{4.4.3}) et de
(\hyperref[II.4.6.2-fr]{4.6.2}) (compte tenu de ce que $f$ est un morphisme
de type fini (\hyperref[II.3.4.1-fr]{3.4.1})).

\section{Morphismes quasi-affines~; morphismes quasi-projectifs. Morphismes
propres~; morphismes projectifs}
\label{section:II.5-fr}

\subsection{Morphismes quasi-affines.}
\label{subsection:II.5.1-fr}

\begin{definition}[5.1.1]
\label{II.5.1.1-fr}
On appelle \emph{schéma quasi-affine} un schéma isomorphe à un sous-schéma
induit sur un ouvert quasi-compact d'un schéma affine. On dit qu'un morphisme
$f:X\to Y$ est \emph{quasi-affine}, ou encore que $X$ (considéré comme
$Y$-préschéma au moyen de $f$) est un \emph{$Y$-schéma quasi-affine}, s'il
existe un recouvrement $(U_\alpha)$ de $Y$ par des ouverts affines tels que
les $f^{-1}(U_\alpha)$ soient des schémas quasi-affines.
\end{definition}

Il est clair qu'un morphisme quasi-affine est \emph{séparé}
(\hyperref[I.5.5.5-fr]{I, 5.5.5} et \hyperref[I.5.5.8-fr]{5.5.8}) et
\emph{quasi-compact} (\hyperref[I.6.6.1-fr]{I, 6.6.1})~; tout morphisme
affine est évidemment quasi-affine.

Rappelons que pour tout préschéma $X$, si on pose
$A=\Gamma(X,\mathcal{O}_X)$, l'homomorphisme identique
$A\to A=\Gamma(X,\mathcal{O}_X)$ définit un morphisme
$X\to\operatorname{Spec}(A)$ dit \emph{canonique}
(\hyperref[I.2.2.4-fr]{I, 2.2.4})~; ce n'est autre d'ailleurs que le morphisme
canonique défini dans (\hyperref[II.4.5.1-fr]{4.5.1}) pour le cas particulier
où $\mathcal{L}=\mathcal{O}_X$, si on se souvient que
$\operatorname{Proj}(A[T])$ s'identifie canoniquement à
$\operatorname{Spec}(A)$ (\hyperref[II.3.1.7-fr]{3.1.7}).

\begin{proposition}[5.1.2]
\label{II.5.1.2-fr}
Soient $X$ un schéma quasi-compact ou un préschéma dont l'espace sous-jacent
est noethérien, $A$ l'anneau $\Gamma(X,\mathcal{O}_X)$. Les conditions
suivantes sont équivalentes~:
\begin{enumerate}
  \item[\rm{a)}] $X$ est un schéma quasi-affine.
  \item[\rm{b)}] Le morphisme canonique
  $u:X\to\operatorname{Spec}(A)$ est une immersion ouverte.
  \item[\rm{b')}] Le morphisme canonique
  $u:X\to\operatorname{Spec}(A)$ est un homéomorphisme de $X$ sur un
  sous-espace de l'espace sous-jacent à $\operatorname{Spec}(A)$.
  \item[\rm{c)}] Le $\mathcal{O}_X$-Module $\mathcal{O}_X$ est très ample
  relativement à $u$ (\hyperref[II.4.4.2-fr]{4.4.2}).
  \item[\rm{c')}] Le $\mathcal{O}_X$-Module $\mathcal{O}_X$ est ample
  (\hyperref[II.4.5.1-fr]{4.5.1}).
  \item[\rm{d)}] Lorsque $f$ parcourt $A$, les $X_f$ forment une base de la
  topologie de $X$.
  \item[\rm{d')}] Lorsque $f$ parcourt $A$, ceux des $X_f$ qui sont affines
  forment un recouvrement de $X$.

